\pagenumbering{roman}

\begin{titlepage}
    \centering
    {\Large Universidad de Guayaquil\par}
    \vspace{1cm}
    {\large Facultad de Jurisprudencia, Ciencias Sociales y Políticas\par}
    \vspace{1cm}
    {\large Carrera de Derecho\par}
    \vspace{2cm}
    {\bfseries\Large Análisis de derecho ambiental ecuatoriano: Revisión crítica de proyectos de ley de apertura al sector minero privado.\par}
    \vspace{2cm}
    {\large Trabajo de titulación\par}
    \vfill
    {\large Autor: Angie Samantha Rivera Vega\par}
    \vspace{0.5cm}
    {\large Guayaquil, Ecuador\par}
    {\large 2026\par}
\end{titlepage}

\chapter*{Dedicatoria}
\addcontentsline{toc}{chapter}{Dedicatoria}

\chapter*{Agradecimiento}
\addcontentsline{toc}{chapter}{Agradecimiento}

\chapter*{Resumen}
\addcontentsline{toc}{chapter}{Resumen}

La apertura del sector minero privado en el Ecuador plantea una tensión jurídica central: determinar si las reformas que la habilitan son compatibles con los principios del derecho ambiental ecuatoriano, los derechos de la naturaleza, la seguridad jurídica y el control estatal. Esta investigación analiza críticamente dicha compatibilidad mediante un enfoque cualitativo, documental y jurídico-crítico que integra análisis normativo, jurisprudencial, doctrinal, institucional e internacional.

El objetivo general consiste en analizar críticamente la compatibilidad de las reformas de apertura minera privada con los principios del derecho ambiental ecuatoriano, los derechos de la naturaleza, la seguridad jurídica y el control estatal aplicable al sector minero en Ecuador.

El análisis se fundamenta en la Constitución de 2008 y el Código Orgánico del Ambiente de 2017, e incorpora jurisprudencia vinculante de la Corte Constitucional sobre derechos de la naturaleza, precaución y consulta ambiental. Se examina la Ley Orgánica para el Fortalecimiento de los Sectores Estratégicos de Minería y Energía, publicada en el Registro Oficial en marzo de 2026; al momento de elaboración de esta tesis, dicha ley acumula demandas de inconstitucionalidad cuyo resultado no ha sido determinado.

La propuesta central es una Matriz de Verificación de Compatibilidad Constitucional y Ambiental con dieciséis criterios de evaluación para proyectos de ley de apertura minera privada. La investigación concluye que la apertura minera privada es constitucionalmente viable únicamente cuando respeta los derechos de la naturaleza, los principios de prevención, precaución y no regresividad, la consulta ambiental, el licenciamiento, la fiscalización estatal efectiva y la reparación integral.

\textbf{Palabras clave:} derecho ambiental; derechos de la naturaleza; minería privada; control estatal; no regresividad; Ecuador.

\chapter*{Abstract}
\addcontentsline{toc}{chapter}{Abstract}

The opening of the private mining sector in Ecuador poses a central legal tension: determining whether the reforms enabling it are compatible with the principles of Ecuadorian environmental law, the rights of nature, legal certainty, and state control. This research critically analyzes such compatibility through a qualitative, documentary, and legal-critical approach that integrates normative, jurisprudential, doctrinal, institutional, and international analysis.

The general objective is to critically analyze the compatibility of private mining sector opening reforms with the principles of Ecuadorian environmental law, the rights of nature, legal certainty, and state control applicable to the mining sector in Ecuador.

The analysis is grounded in the 2008 Constitution and the 2017 Organic Environmental Code, incorporating binding jurisprudence from the Constitutional Court on the rights of nature, the precautionary principle, and environmental consultation. The Organic Law for the Strengthening of the Strategic Mining and Energy Sectors, published in the Official Registry in March 2026, is examined; at the time of elaboration of this thesis, said law accumulates unconstitutionality claims whose outcome has not been determined.

The central proposal is a Constitutional and Environmental Compatibility Verification Matrix with sixteen evaluation criteria for private mining sector opening bills. The research concludes that such opening is constitutionally viable only when it respects the rights of nature, the principles of prevention, precaution, and non-regression, environmental consultation, licensing, effective state oversight, and full reparation.

\textbf{Keywords:} environmental law; rights of nature; private mining; state control; non-regression; Ecuador.

\clearpage
\pagenumbering{arabic}